\subsection{Actual finite curve normalization and original ideal pullback}
\label{subsec:finite-curve-normalization}

\begin{definition}[Actual normalization ring in the original fraction field]
\label{algebra:actual-curve-normalization-ring}
\lean{EnvelopeProperty.CurveNormalization.Ring,
EnvelopeProperty.CurveNormalization.inclusion,
EnvelopeProperty.CurveNormalization.inclusion_coe,
EnvelopeProperty.CurveNormalization.inclusion_injective,
EnvelopeProperty.CurveNormalization.isFractionRing_ring,
EnvelopeProperty.CurveNormalization.isIntegrallyClosed_ring}
\leanok
Let $B$ be an actual commutative domain. Its normalization ring is the
actual integral closure $C$ of $B$ in its original fraction field
$L=\operatorname{Frac}B$. The original inclusion $B\to C$ retains every
original fraction-field element and is injective. The entire fraction
field of $C$ is this same $L$, and $C$ is integrally closed.
The underlying integral-closure construction also exists in the fraction
ring of an arbitrary commutative ring; the domain normality and
fraction-field assertions here retain their explicit domain hypothesis.
\end{definition}
\begin{proof}
Take the actual subtype of integral elements in the original fraction
ring and its literal original algebra map. The entire field image of
each included element is the original map $B\to L$. This map is
injective, so subtype equality proves actual inclusion injectivity.
The genuine integral-closure localization theorem identifies the
original field $L$ as the fraction field of $C$, taking the finite
self-extension $L/L$. Integral closure in a field is integrally closed,
by the actual idempotence and integral transitivity theorem. No
separability or finite-generation hypothesis is needed for these
whole-ring assertions.
\end{proof}

\begin{theorem}[Actual finite, Dedekind and free curve normalization]
\label{algebra:finite-free-curve-normalization}
\lean{EnvelopeProperty.CurveNormalization.coefficientMap,
EnvelopeProperty.CurveNormalization.coefficientMap_injective,
EnvelopeProperty.CurveNormalization.finite_inclusion,
EnvelopeProperty.CurveNormalization.isNoetherian_ring,
EnvelopeProperty.CurveNormalization.finite_coefficientMap,
EnvelopeProperty.CurveNormalization.isDedekindDomain_ring,
EnvelopeProperty.CurveNormalization.free_coefficientMap,
EnvelopeProperty.CurveNormalization.flat_coefficientMap}
\leanok
\uses{algebra:actual-curve-normalization-ring,thm:bj-normalization-finite}
Suppose $B$ is a finite-type domain over a perfect field $F$. Its
actual normalization inclusion is finite and $C$ is Noetherian.
If an actual $A$-algebra structure makes $B$ finite over $A$, the literal
composite $A\to B\to C$ is finite. If $A$ has dimension at most one,
$C$ is a genuine Dedekind domain. If $A$ is a principal ideal domain
and the actual original $A\to B$ is injective, $C$ is a free $A$-module
for that precise composite map, and the map is flat. Arbitrary
characteristic is retained; no separability hypothesis is imposed.
\end{theorem}
\begin{proof}
Apply the checked perfect-field finite-normalization theorem to the
literal integral closure of $B$ in its fraction field. This theorem's
actual Frobenius finite-module descent covers inseparable cases as
well as the separable trace-form case. Thus $B\to C$ is finite and
$C$ is Noetherian. Composing with the original finite $A\to B$ proves
finiteness of the whole original composite, rather than choosing a
new coefficient action.

This actual finite composite is integral. The integral dimension
comparison gives dimension at most one for $C$ when the same holds
for $A$. Together with actual Noetherianity, domain structure and the
proved integrally closedness, this is precisely the genuine Dedekind
property. If $A\to B$ is injective, compose with the proved normalization
inclusion to get injectivity of the entire map $A\to C$. Since both
rings are domains, the actual resulting coefficient module is
torsion-free. A genuine finite torsion-free module over a principal
ideal domain is free, and a true free module is flat. The exported
free-module assertion explicitly uses the algebra action of this
literal composite homomorphism. The original base-map injectivity
is a primitive algebraic hypothesis, supplied independently by the
actual line-component dimension and domination producers.
\end{proof}

\begin{theorem}[Actual finite normalization and coefficient-curve scheme maps]
\label{scheme:actual-finite-flat-curve-normalization}
\lean{EnvelopeProperty.CurveNormalization.scheme,
EnvelopeProperty.CurveNormalization.projection,
EnvelopeProperty.CurveNormalization.curveProjection,
EnvelopeProperty.CurveNormalization.projection_coefficientMap,
EnvelopeProperty.CurveNormalization.projection_isFinite,
EnvelopeProperty.CurveNormalization.curveProjection_isFinite,
EnvelopeProperty.CurveNormalization.curveProjection_flat,
EnvelopeProperty.CurveNormalization.projection_surjective}
\leanok
\uses{algebra:finite-free-curve-normalization,algebra:actual-curve-normalization-ring,
scheme:integral-native-proper-map,scheme:integral-native-image}
The actual affine normalization scheme is $N=\operatorname{Spec}C$
with its literal original projection $N\to\operatorname{Spec}B$.
It is finite under the preceding perfect-field finite-type hypotheses
and is surjective on actual scheme points. Its entire composite to
$\operatorname{Spec}A$ is the actual composite coefficient scheme map,
finite whenever $B$ is finite over $A$ and flat under the preceding
actual principal-domain and injectivity hypotheses.
\end{theorem}
\begin{proof}
Apply the genuine spectrum functor to the literal original ring maps.
Its contravariant composition law gives the whole original coefficient
triangle. The actual affine finiteness criterion transfers each proved
finite ring map to the corresponding true scheme morphism; the affine
flatness criterion transfers the genuine flat ring map likewise.
The normalization inclusion is integral by the actual defining integral
elements, and it is injective by the preceding original fraction-field
identity. True integral lying-over therefore proves surjectivity on
the entire prime spectrum, which is the actual underlying scheme map.

The independently proved general native integral-map producers apply
to this real integral projection: its native map is proper, and the
actual scheme surjectivity yields whole native surjectivity, also on
an exact common prescribed-base locus. Thus these existing producers
supply genuine native lifts without a bound on a prescribed field
valuation. No particular normalization branch over a given point,
universal injectivity or homeomorphism of a general normalization
projection is asserted. Smoothness of the normalized curve is not
inferred from the Dedekind assertion without a genuine regular-to-smooth
producer.
\end{proof}

\begin{theorem}[Original affine and stalk ideal pullback through normalization]
\label{scheme:normalization-original-ideal-pullback}
\lean{EnvelopeProperty.CurveNormalization.idealSheaf,
EnvelopeProperty.CurveNormalization.idealSheaf_projection_comap,
EnvelopeProperty.CurveNormalization.stalkIdeal_projection_comap}
\leanok
\uses{scheme:actual-finite-flat-curve-normalization,scheme:genuine-ideal-pullback-stalk}
For every original ideal $J\subseteq B$, its actual affine ideal-sheaf
datum pulls back through the literal normalization projection to the
actual ideal datum of the entire image $JC\subseteq C$. For every
actual ideal-sheaf datum $I$ on $\operatorname{Spec}B$ and every actual
normalization point, the intrinsic stalk ideal of its geometric
pullback is the image of the original whole stalk ideal under the
true normalization stalk map. No radical replacement is made.
\end{theorem}
\begin{proof}
Construct the original affine ideal datum by mapping $J$ through the
actual global-section isomorphism of its spectrum and applying the
genuine affine ideal-sheaf construction. The independently proved
geometric pullback theorem makes the actual affine section ideal of
its pullback the image of this whole original section ideal. Genuine
naturality of the spectrum/global-section isomorphism identifies the
entire section map with the literal original inclusion $B\to C$.
Ideal-image composition therefore identifies this whole ideal with
the original image $JC$ under the true source global-section
isomorphism. The genuine affine ideal-sheaf equivalence identifies
these entire geometric ideal data. In particular original powers and
nonradical ideals are retained.

For stalks, apply the actual all-scheme geometric stalk-image theorem
to this literal normalization morphism. It is proved from true affine
refinement, the geometric ideal pullback and whole germ/stalk
naturality; no finite generation, properness or stalk-image certificate
is needed. These statements concern actual algebraic sheaf and stalk
pullback, not a whole-field ideal supremum substituted for a model norm.
\end{proof}

\begin{theorem}[Separate given bounded-model norm consistency]
\label{scheme:normalization-given-bounded-model-functions}
\lean{EnvelopeProperty.CurveNormalization.modelIdealNorm_projection_comap,
EnvelopeProperty.CurveNormalization.modelIdealLog_projection_comap}
\leanok
\uses{scheme:actual-finite-flat-curve-normalization,
scheme:normalization-original-ideal-pullback,scheme:proper-model-ideal-pullback}
Separately, suppose an actual given model structure
$f:\operatorname{Spec}B\to\operatorname{Spec}R$ is proper, and a
prescribed real seminorm $v$ on $R$ satisfies $v(r)\le1$ for every
model coefficient. For every original ideal-sheaf datum $I$ and every
native normalization-model point above this exact $v$, its actual
geometric pullback has intrinsic model norm and logarithm equal to
those of $I$ at the whole original native image point.
\end{theorem}
\begin{proof}
The actual finite normalization projection is proper. Compose it
with the given proper model structure to derive properness of the
source model over the same actual coefficient ring $R$. No properness
of this normalized model is introduced as an independent certificate.
The explicit bound on $v$ supplies genuine proper centers on both
models. Apply the proved whole geometric ideal-pullback model-norm
identity to the literal normalization projection and its whole
coefficient triangle. Its actual center-map comparison and bounded
center-stalk values compute the original intrinsic ideals, proving
the entire norm identity. Taking the genuine extended logarithm gives
the second identity, including zero and negative infinity.

This is a corollary for a given bounded proper affine model; it is
not the paper's field-curve model construction. For a nontrivially
valued $F$, such an affine proper bounded model with $B$ containing
$F$ does not generally supply native points with that prescribed
field valuation. These statements neither construct those models
or points nor replace their model ideal functions by intrinsic
\texttt{idealNorm} over all coefficients of the valued field, where
extension can make a nonzero ideal a unit or an unbounded whole ideal.
The field-curve scope above is the actual finite normalization map
and original geometric section/stalk pullback. Actual proper-model
existence for that curve, smoothness, analytic subharmonic comparison
and continuation remain separate producers.
\end{proof}
